\section{Introduction}

Point clouds have become a common representation for scanned buildings, terrain, industrial sites, robotics environments, and reconstructed cultural heritage. Their appeal is directness: a point set can preserve fine geometric detail without requiring a watertight surface model. The same directness creates a practical systems problem. A cloud with millions or hundreds of millions of samples is only useful if spatial queries can be answered quickly enough for the surrounding application, and the best index for those queries is often not obvious before the workload is known.

The standard answer is to choose one spatial data structure and tune it by hand. Quadtrees and kd-trees are classical examples of this design space~\cite{finkel1974quadtrees,bentley1975kdtree}; related spatial-index families such as R-trees and bounding-volume hierarchies expose similar trade-offs between subdivision, bounds, locality, and update cost~\cite{guttman1984rtree,karras2012lbvh}. Quadtrees and octrees are simple and predictable for axis-aligned subdivision. kd-trees adapt naturally to anisotropic distributions. Bounding-volume hierarchies can be effective when tight bounds matter. Uniform and hierarchical grids are attractive when constant-time cell lookup outweighs their sensitivity to density variation. Each of these structures has a region where it behaves well, but point-cloud workloads do not stay inside one region. A facade scan may be almost planar in one area and cluttered in another. A terrestrial laser scan may combine ground, walls, vegetation, and empty space. A visualization tool may issue small neighborhood queries, while an analytics task may issue large range filters over the same data.

Recent work confirms that this is not only a matter of choosing better defaults. Acceleration-structure autotuning has shown that BVH and kd-tree parameters can be searched for ray-tracing workloads~\cite{herveau2023autotuning}. Point-cloud systems have tuned octree and R*-tree thresholds, built nested octrees for BIM-to-point-cloud change detection, and proposed user-configurable nested spatial structures for LiDAR data~\cite{wang2021hybridpointcloud,park2021nestedoctree,ogayar2023nestedlidar}. In parallel, learned and workload-aware indexing systems optimize layouts or hierarchy construction from query distributions~\cite{sheng2023wisk,pai2023wazi,wang2025litune}. These results support a broader premise: index performance depends jointly on data distribution, workload, hardware, and the chosen structural parameters. They also expose a gap. For three-dimensional point-cloud and mesh-style workloads, the literature is stronger on individual parameter tuning and hand-designed hybrids than on automatic search over heterogeneous, deeply composable index schemas.

This paper treats point-cloud index selection as physical design. It is also a form of data-structure meta-optimization: the system searches over designs for the data and workload at hand, rather than assuming that one fixed primitive should dominate. Recent work on learned indexes, the Data Calculator, and design continuums argues that data structures should be understood less as a fixed catalogue and more as a parameterized design space whose best point depends on data, workload, hardware, and constraints~\cite{kraska2018learned,idreos2018datacalculator,idreos2019learningkv}. We adopt that view for spatial point-cloud indexes, but keep the final decision measurement-driven. Given a point set and a query profile, the system searches over a space of replayable schemas rather than over one primitive at a time. A schema is an ordered description of one or more spatial-index levels: for example, an octree followed by a kd-tree inside selected leaves, or a quadtree that activates deeper three-dimensional subdivision only in nodes whose local geometry justifies it. The purpose is not to replace the well-known primitives. The purpose is to make them composable, measurable, and selectable under the workload that will actually be run.

\begin{figure}[t]
  \centering
  \figureplaceholder{0.26\textheight}{Pipeline diagram: point cloud and workload profile enter feature extraction; generated and configured schemas form a candidate pool; candidates are measured on CPU or CUDA backends; the output is a selected replayable schema, raw measurement table, and optional selector-training data.}
  \caption{Overview of the proposed workload-aware schema synthesis workflow. The final figure should show the optimizer as a measured loop rather than a purely analytical predictor: candidate schemas are built, queried, scored, audited for provenance, and exported for replay.}
  \label{fig:pipeline}
\end{figure}

Our approach is deliberately empirical. Learned or analytical cost models are valuable for exploring large design spaces, but they become brittle when node occupancy, branch conditions, memory layout, backend differences, and query mix interact. Instead, we treat a candidate schema as a hypothesis. The framework builds the index, runs a deterministic query workload, records latency and structural counters, and uses those measurements to select or improve the next candidates. This makes the optimization loop closer to autotuning than to a closed-form data-structure analysis. It also keeps the output inspectable: the selected schema is a JSON artifact that can be replayed, compared against single-primitive baselines, and used as a label for a later learned selector.

The resulting system, MultiDataStructure, is aimed at large point-cloud indexing rather than relational query processing. It does not optimize SQL plans, and it does not assume a fixed database storage layer. Its object of optimization is the spatial index itself: the primitive sequence, subdivision depth, leaf capacity, split policy, conditional gates, and optional local leaf indexes. The maintained query interface supports exact axis-aligned range, count-range, radius, and k-nearest-neighbor queries over a loaded point cloud. Candidate schemas can be evaluated on a CPU reference implementation and, where supported, on CUDA builders for GPU-resident variants.

The design follows three principles. First, all comparisons should be workload aware. A schema is not declared good in isolation; it is measured under a query distribution with explicit range scales, radius scales, query counts, seeds, and score weights. Second, all comparisons should keep provenance. The framework records the backend, CUDA builder, query count, score stage, score weights, support status, and whether a row is a final latency measurement or a proxy-stage measurement. Third, the optimizer should be allowed to choose a simple answer. If a tuned single octree is best for a given cloud and workload, the search should return that single primitive rather than forcing a nested design.

This draft focuses on the method. The contributions of the completed paper will be:
\begin{itemize}[leftmargin=*]
  \item a schema representation that expresses single, nested, conditional, and adaptive point-cloud indexes using the same replayable format;
  \item a measured schema-search pipeline that combines deterministic workload generation, structural diagnostics, staged candidate evaluation, and baseline controls;
  \item a spatial instance of data-structure meta-optimization, where mature index primitives are recombined and tuned under measured point-cloud workloads rather than replaced by a single monolithic learned model;
  \item a feature protocol for connecting point-cloud geometry, workload intent, and schema composition to later learned ranking;
  \item CPU and CUDA evaluation paths that expose both algorithmic behavior and backend-specific support limitations instead of hiding fallback cases.
\end{itemize}

The rest of this manuscript is organized as follows. Section~\ref{sec:methodology} defines the point-cloud and workload model, the schema space, the scoring objective, the staged search procedure, and the measurement protocol. The evaluation section will later report controlled comparisons against single-primitive and hand-designed nested baselines. The discussion will cover failure modes, GPU support gaps, and the boundary between per-cloud tuning and general learned selection.
